% !TEX program = pdflatex
\documentclass[11pt,a4paper]{article}
\usepackage[margin=25mm]{geometry}
\usepackage[T1]{fontenc}
\usepackage[utf8]{inputenc}
\usepackage{lmodern}
\usepackage{amsmath,amssymb,amsthm,mathtools}
\usepackage{booktabs,array,tabularx}
\usepackage{microtype}
\usepackage[hidelinks,unicode]{hyperref}
\hypersetup{pdftitle={Identités continues, formalisation et budgets explicites pour Goldbach binaire -- Synthèse V3},pdfauthor={Durand Serge},pdfsubject={Synthèse de recherche et frontière de vérification},pdfkeywords={Goldbach binaire, von Mangoldt, Mellin, projection, Lean 4}}
\setlength{\emergencystretch}{3em}
\setlength{\parskip}{3pt}
\setlength{\tabcolsep}{4pt}
\renewcommand{\arraystretch}{1.14}
\newcolumntype{Y}{>{\raggedright\arraybackslash}X}
\newtheorem{theorem}{Théorème}[section]
\newtheorem{proposition}[theorem]{Proposition}
\newtheorem{definition}[theorem]{Définition}
\theoremstyle{remark}
\newtheorem{remark}[theorem]{Remarque}
\newcommand{\Lam}{\Lambda}
\newcommand{\Sing}{\mathfrak S}
\newcommand{\C}{\mathbb C}
\newcommand{\R}{\mathbb R}
\newcommand{\N}{\mathbb N}
\newcommand{\Z}{\mathbb Z}
\newcommand{\paid}{\textbf{Preuve Lean compilée.}\space}
\newcommand{\source}{\textbf{SOURCE non compilée.}\space}
\newcommand{\paper}{\textbf{PAPER : argument écrit.}\space}
\DeclareMathOperator{\Arg}{Arg}
\DeclareMathOperator{\Log}{Log}
\DeclareMathOperator{\sech}{sech}
\title{Identités continues, formalisation et budgets explicites\\
pour Goldbach binaire\\[5pt]\large Synthèse V3 et clôture d'étape}
\author{Durand Serge\\\small Chercheur indépendant\\\small\texttt{seisner1967@gmail.com}}
\date{4 octobre 2026}
\begin{document}
\maketitle
\begin{abstract}
Cette synthèse poursuit la monographie V2 en conservant son cadre bilantiel et
ses notations. La recherche de la boucle 22 transpose le couplage additif dans
des traces continues, des intégrales de Mellin et des projections de caractères.
L'inventaire arrêté après le lot indépendant 35 contient 88 modules et
1\,488 déclarations auxiliaires, définitions incluses. Les acquis nouveaux
comprennent l'extraction exacte du coefficient de la vraie fonction de von
Mangoldt sur le cercle, sa projection discrète sans alias, une quantification
canonique des logarithmes à l'échelle $2^{58}$, les cinq racines de corps finis
du contrat, et une représentation de Mellin avec échange somme--intégrale
effectivement justifié. Les enveloppes de troncature pondérées, leur géométrie
uniforme et leur raccord complet au coefficient demeurent des sources de
preuve non compilées. Un test rationnel exact valide un majorant scalaire
inférieur à $10^{-6}$ au point $N=10^8$ ; il n'évalue pas le coefficient.
Le contrôleur natif indépendant s'arrête sans verdict mathématique à sa limite
de temps. La borne cible pour $D_N$, la contribution signée globale et la
conjecture de Goldbach binaire restent ouvertes.
\end{abstract}
\noindent\textbf{Mots-clés :} Goldbach binaire ; von Mangoldt ; Mellin ;
projection continue ; erreurs explicites ; Lean 4.
\par\noindent\textbf{Classification :} 11P32.
\par\noindent\textbf{Portée.} Ce texte est une clôture de recherche et un
contrat de reprise. Il ne contient aucune annonce de preuve de Goldbach ni de
franchissement du mur de la parité. Les statuts se rapportent aux reçus
conservés, et non au nombre d'identités que l'on pourrait encore formaliser.
\par\noindent{\small Le manuscrit a été préparé avec assistance d'agents de
recherche et d'outils de rédaction. L'assistance ne constitue pas une
validation scientifique externe. Les preuves compilées sont identifiées par
leurs sources et les audits indépendants de leurs axiomes.}

\clearpage
\section{Cadre conservé et frontière de vérification}\label{sec:frame}
Pour un entier pair $N$, on conserve les notations
\[
u=\log N,\qquad L=1+u,\qquad \ell=\log u,
\]
avec logarithmes naturels. Le seuil du cadre analytique fixé est
\begin{equation}\label{eq:onset}
u\ge 10^{24}.
\end{equation}
Le point $N=10^8$ sert de boussole finie ; il se situe hors de ce régime.
Une vérification à ce point ne couvre ni le seuil \eqref{eq:onset}, ni tous les
entiers pairs en dessous de celui-ci. La présente étape conserve les acquis
du cadre et examine leur raccord à de nouveaux objets continus.

La référence du bilan reste
\begin{equation}\label{eq:reference}
R_{\mathrm{ref}}=\Sing(N)(N-F_N),\quad
\Sing(N)=2C_2\prod_{\substack{p\mid N\\p>2}}\frac{p-1}{p-2},\quad
C_2=\prod_{p>2}\left(1-\frac1{(p-1)^2}\right).
\end{equation}
Les dépendances analytiques antérieures de $F_N$ et de la référence gardent
leur statut dans la monographie. Cette synthèse ne transforme pas une entrée
analytique rapportée en un nouveau théorème Lean autoportant. La cible de la
recherche demeure
\begin{equation}\label{eq:target}
D_N\le \tau_N:=\frac{N}{256u\ell}.
\end{equation}
Les coupures du profil conservé sont
\[
\alpha=\lceil N^{1/4}\rceil,\quad Q=\left\lfloor\frac{N-1}{\alpha}\right\rfloor,
\quad a_{\mathrm{cut}}=\lceil N^{7/16}\rceil,
\quad M_{\mathrm{cut}}=\lceil N^{3/4}\rceil.
\]
Le paramètre thermique $a>0$ introduit plus loin est distinct de
$a_{\mathrm{cut}}$. Tous les fronts, bandes, coins, unités et restrictions
du support bilantiel restent à leurs places.

\subsection*{Quatre niveaux de preuve}
\begin{tabularx}{\textwidth}{@{}p{29mm}Y@{}}
\toprule
Niveau & Sens exact dans cette synthèse\\
\midrule
Lean compilé & Source gelée, compilation indépendante réussie, catalogue
exact et audit des axiomes ; crédit au module entier.\\
SOURCE & Texte de preuve Lean rédigé et éventuellement relu ; élaboration
et compilation encore absentes ou échouées.\\
PAPER & Identité ou estimation dérivée dans une note ; ses obligations de
formalisation et d'évaluation demeurent explicites.\\
Calcul observé & Programme effectivement lancé sous un contrat précis ;
son résultat n'est crédité qu'à la portée de ce contrat.\\
\bottomrule
\end{tabularx}
Les 1\,488 déclarations comptées incluent les définitions. Elles ne représentent
pas 1\,488 théorèmes de Goldbach. Les axiomes usuels apparaissant dans les
audits comprennent \texttt{propext}, \texttt{Classical.choice} et
\texttt{Quot.sound}. Aucun \texttt{sorryAx} n'est accepté dans un module crédité.
Les sorties de récupération après une élaboration échouée restent des faits
d'échec technique, sans crédit de preuve.

\clearpage
\section{La vraie trace thermique et le coefficient additif}\label{sec:heat}
On utilise la fonction arithmétique réelle
\[
\Lam(0)=\Lam(1)=0,\qquad
\Lam(n)=\begin{cases}\log p,&n=p^k,\ p\text{ premier},\ k\ge1,\\0,&\text{sinon}.
\end{cases}
\]
Toutes les puissances premières sont conservées. Pour $a>0$ et $\theta\in\R$,
\begin{equation}\label{eq:heat}
T_a(\theta)=\sum_{n\ge0}\Lam(n)e^{-an}e^{in\theta},\qquad
C_N=\sum_{n=0}^{N}\Lam(n)\Lam(N-n).
\end{equation}
L'usage d'une fonction de trace abstraite n'est pas nécessaire : la série
ci-dessus est l'objet effectivement formalisé.

Posons $r=e^{-a}\in(0,1)$ et
\begin{align}
U(a)&=\frac{r}{(1-r)^2},\label{eq:U}\\
B(a,M)&=r^{M+1}\left(\frac{M+1}{1-r}+U(a)\right).
\label{eq:heat-tail}
\end{align}
\paid Le lot 11 construit directement $0\le\Lam(n)\le n$, la
sommabilité des termes, la continuité de $T_a$, et les bornes uniformes
\[
|T_a(\theta)|\le U(a),\qquad
|T_a(\theta)-T_{a,M}(\theta)|\le B(a,M),
\quad T_{a,M}(\theta)=\sum_{n=0}^{M}\Lam(n)r^ne^{in\theta}.
\]
Il prouve aussi la continuité des enveloppes en $a>0$, et l'erreur entre
les projections intégrales complète et partielle, majorée par
$2e^{aN}U(a)B(a,M)$. Le majorant géométrique est construit ; il n'est pas
fourni comme une hypothèse de précision.

\begin{theorem}[Projection continue exacte ; lot 13]\label{thm:circle}
Pour tout $a>0$ et tout $N\in\N$,
\begin{equation}\label{eq:circle}
C_N=\frac{e^{aN}}{2\pi}\int_0^{2\pi}
T_a(\theta)\,\overline{T_a(-\theta)}\,e^{-iN\theta}\,d\theta
=\frac{e^{aN}}{2\pi}\int_0^{2\pi}T_a(\theta)^2e^{-iN\theta}\,d\theta.
\end{equation}
\end{theorem}
\paid La preuve extrait d'abord les caractères sur un rectangle fini,
paie l'orthogonalité intégrale et le facteur de normalisation, puis passe à
la limite avec l'enveloppe \eqref{eq:heat-tail}. Dès $M\ge N$, la projection
intégrale du polynôme $T_{a,M}$ vaut déjà exactement $C_N$.

La conjugaison porte sur la phase réfléchie :
$\overline{T_a(-\theta)}=T_a(\theta)$. Le produit $|T_a(\theta)|^2$
sélectionnerait des différences de fréquences ; il ne fournit pas le
coefficient additif recherché. La forme centrée sur $[-\pi,\pi]$ repose sur
la périodicité de la vraie série. Son assemblage avec la troncature de Mellin
est rédigé séparément, au statut SOURCE décrit en section~\ref{sec:radius}.

\clearpage
\section{Projecteur fini sans alias et projection de corps fini}\label{sec:finite}
L'identité continue possède une variante entièrement finie. Elle sert de
contrat exact à un calcul éventuel, sans remplacer la vraie fonction $\Lam$.
Pour $K\ge1$, posons $\theta_j=2\pi j/K$.

\begin{theorem}[Orthogonalité discrète et extraction ; lot 15]
Pour $N\le M$ et
\begin{equation}\label{eq:alias}
K>\max(N,2M-N),
\end{equation}
on a, pour tout réel $a$,
\begin{equation}\label{eq:discrete}
C_N=\frac{e^{aN}}{K}\sum_{j=0}^{K-1}
T_{a,M}(\theta_j)^2e^{-iN\theta_j}.
\end{equation}
\end{theorem}
\paid Le lot 15 prouve l'orthogonalité des véritables caractères
exponentiels, puis leur expansion et la sélection de l'antidiagonale.
La condition \eqref{eq:alias} interdit les fréquences parasites
$m+n-N\in K\Z$ avec $0\le m,n\le M$. Elle est une garde mathématique,
et non une tolérance expérimentale. Le choix $a=0$ est licite dans ce
polynôme fini ; aucune trace infinie $T_0$ n'est introduite.

Une projection analogue existe dans un corps fini. Pour un nombre premier
$p$, un élément $\omega\in\mathbb F_p$ d'ordre exactement $K$ et
$K\ne0$ dans $\mathbb F_p$, l'orthogonalité donne
\[
\frac1K\sum_{j=0}^{K-1}\omega^{jk}
=\begin{cases}1,&K\mid k,\\0,&K\nmid k.\end{cases}
\]
\paid Le lot 19 prouve cette identité depuis la racine primitive ;
l'orthogonalité finale n'est pas une prémisse. Pour une suite finie $A_n$ et
\eqref{eq:alias}, sa projection sélectionne
$\sum_{n=0}^N A_nA_{N-n}$ dans $\mathbb F_p$.

La représentation par un cercle continu et cette projection de caractères
ont un objet commun : le coefficient additif exact. Une identité de
projection ne dit pas que ce coefficient est assez grand. Elle n'établit
pas non plus la conformité d'un programme de transformée rapide à la somme
finie du théorème. Cette conformité exige les invariants de ses étages, les
normalisations, les bornes des mots machine et la reconstruction entière.

Le contrat final retient
\begin{equation}\label{eq:fixed}
N=M=100\,000\,000,\qquad K=2^{27}=134\,217\,728,
\qquad S=2^{58}=288\,230\,376\,151\,711\,744.
\end{equation}
Ici $K>N$ satisfait \eqref{eq:alias}. Les preuves finies demeurent valides
indépendamment de l'achèvement d'un calcul à cette taille. Le contrat de
validation doit conserver les coefficients issus de toutes les puissances
premières, y compris celles de degré au moins deux.

\clearpage
\section{Logarithmes canoniques et précision effectivement construite}\label{sec:quant}
La quantification du lot 17 ne suppose pas une précision fournie par un
oracle. Pour $2\le p\le10^8$, posons
\[
k=\lfloor\log_2p\rfloor,\qquad
z=\frac{p-2^k}{p+2^k},\qquad 0\le z<\frac13,
\]
et définissons la somme rationnelle
\begin{equation}\label{eq:log32}
L_{32}(z)=2\sum_{j=0}^{31}\frac{z^{2j+1}}{2j+1},\qquad
\rho_{32}=\frac{9}{4\cdot65\cdot3^{65}}.
\end{equation}
La représentation réelle de $\log(1+z)-\log(1-z)$ par cette série,
sa queue géométrique et la décomposition par $2^k$ donnent l'intervalle
\begin{equation}\label{eq:logbox}
b_p=kL_{32}(1/3)+L_{32}(z),\qquad
\log p\in[b_p,b_p+(k+1)\rho_{32}].
\end{equation}
\paid La largeur satisfait, sur le domaine fixé,
\[
(k+1)\rho_{32}\le\frac{243}{260\cdot3^{65}}<2^{-96}.
\]
Le point canonique $A_p$ est obtenu en multipliant le milieu de
\eqref{eq:logbox} par $S$, en arrondissant à l'entier le plus proche avec
égalité départagée vers l'entier pair, puis en le bornant dans $[0,32S]$.
Le constructeur utilise quotient et reste rationnels exacts. Il n'est pas
identifié sans preuve à une fonction d'arrondi différente.

Pour tous les entiers, définissons
\[
A_n=\begin{cases}
A_{\min\mathrm{Fac}(n)},& n\text{ est une puissance première},\\
0,&\text{sinon},
\end{cases}\qquad q_n=A_n/S.
\]
Les cas $0$ et $1$ sont traités séparément. L'emploi de la définition directe
de $\Lam$ et de son facteur premier minimal conserve chaque $p^k$.
\paid La preuve établit, pour $0\le n\le10^8$,
\begin{equation}\label{eq:q-error}
0\le\Lam(n),q_n\le32,\qquad |\Lam(n)-q_n|\le 1/S.
\end{equation}
La précision provient de \eqref{eq:log32}--\eqref{eq:logbox} et de l'arrondi
construit ; elle n'apparaît pas comme hypothèse libre du théorème.

Posons $I_N=\sum_{n=0}^{N}A_nA_{N-n}\in\Z_{\ge0}$. Le même lot prouve
\begin{equation}\label{eq:E-log}
\left|C_N-\frac{I_N}{S^2}\right|\le E_{\log}(N,S),\qquad
E_{\log}(N,s)=(N+1)\left(\frac{64}{s}+\frac1{s^2}\right),\quad s>0.
\end{equation}
L'enveloppe est continue en $s>0$. La garde entière fixée
\begin{equation}\label{eq:q-guard}
2(N+1)(64S+1)10^6\le S^2
\end{equation}
est compilée, donc $2E_{\log}(10^8,2^{58})\le10^{-6}$. Ce résultat paie
la quantification réelle, sans calculer $I_N$ ni $C_N$.

\clearpage
\section{Cinq modules premiers, résidus exacts et dette native}\label{sec:crt}
Le contrat fixe les cinq paires suivantes ; leur primalité et l'ordre
des racines ont été compilés indépendamment au lot 22.
\begin{center}
\begin{tabular}{@{}rrrr@{}}
\toprule
$i$ & $p_i$ & $(p_i-1)/K$ & $g_i$\\
\midrule
1 & 2\,013\,265\,921 & 15 & 31\\
2 & 2\,281\,701\,377 & 17 & 3\\
3 & 3\,221\,225\,473 & 24 & 5\\
4 & 3\,489\,660\,929 & 26 & 3\\
5 & 3\,892\,314\,113 & 29 & 3\\
\bottomrule
\end{tabular}
\end{center}
Dans $\mathbb F_{p_i}$, on pose
$\omega_i=g_i^{(p_i-1)/K}$. La preuve établit l'ordre $K$, à partir
des certificats de primalité et des contrôles des puissances nécessaires.

\begin{theorem}[Projection canonique concrète ; lot 22]
Avec les paramètres \eqref{eq:fixed} et les entiers canoniques $A_n$,
\begin{equation}\label{eq:ntt-projection}
F_i(j)=\sum_{n=0}^{N}A_n\omega_i^{jn},\qquad
r_i=K^{-1}\sum_{j=0}^{K-1}F_i(j)^2\omega_i^{-Nj}
=I_N\pmod {p_i}.
\end{equation}
\end{theorem}
\paid Le théorème instancie la projection de corps fini avec les
vraies racines du tableau. Il ne prend ni le résidu final ni la précision
des coefficients en prémisse.

\paper Pour $\Pi=\prod_{i=1}^5p_i$, les bornes du contrat sont
\begin{equation}\label{eq:crt-size}
0\le I_N\le(N+1)(32S)^2<2^{153}<\Pi.
\end{equation}
Elles fournissent le domaine de reconstruction unique par restes chinois.
Une reconstruction $c\in[0,\Pi)$ ayant les cinq bons résidus doit donc être
$I_N$, sous le raccord exact au constructeur et aux projections.
La preuve du comportement du code CRT, de ses mots machine et de la
transformée rapide reste distincte de \eqref{eq:ntt-projection}.

Le producteur et le checker natifs ont tous deux été construits au
BUILD04 : deux compilations terminées avec code zéro. Cela établit
l'existence de deux images exécutables gelées, sans prouver leur raffinement
par rapport aux fonctions Lean. Le checker doit encore certifier tout le
catalogue jusqu'à $N$, la classification des puissances premières, l'égalité
de chaque point enregistré au constructeur canonique $A32$, les résidus,
la reconstruction et les gardes entières. Un contrôle auxiliaire par une
série $B40$ n'est pas le constructeur primaire $A32$ ; son code et sa
précision réelle n'ont pas de preuve Lean complète dans cette étape.

La fenêtre de ressources est une contrainte opérationnelle. La confiance
déclarée en Windows, GCC et les images natives n'est ni une fermeture
universelle du chargeur ni une observation de toutes les dépendances
dynamiques. Aucune de ces hypothèses opérationnelles ne remplace un
théorème de raffinement du programme.

\clearpage
\section{Mellin : identité exacte avec la vraie série de von Mangoldt}\label{sec:mellin}
Soit $w\in\C$ tel que $\Re w>0$. Le logarithme principal détermine
$w^{-s}=\exp(-s\Log w)$. Posons
\begin{equation}\label{eq:mellin-objects}
s(t)=2+it,\quad
D(t)=\sum_{n\ge2}\Lam(n)n^{-2-it},\quad
P(w)=\sum_{n\ge0}\Lam(n)e^{-nw},\quad
K(w,t)=\Gamma(2+it)w^{-2-it}.
\end{equation}
Les termes de rang zéro et un de $D$ sont nuls dans la définition formelle. Pour
$n\ge1$, les puissances utilisent le logarithme réel positif de $n$.
En particulier $P(a-i\theta)=T_a(\theta)$.

\begin{theorem}[Inversion et échange infini ; lot 35]\label{thm:mellin}
Pour $\Re w>0$, la fonction $t\mapsto D(t)K(w,t)$ est intégrable et
\begin{equation}\label{eq:mellin}
P(w)=\frac1{2\pi}\int_{\R}D(t)\Gamma(2+it)w^{-2-it}\,dt.
\end{equation}
De plus $D$ est continue et $|D(t)|\le6$ pour tout $t\in\R$.
\end{theorem}
\paid Le lot 35 contient 30 déclarations : 24 théorèmes et 6
définitions. Sa compilation indépendante et l'audit exact de ses axiomes
ont réussi. L'échange ne repose pas sur une hypothèse finale d'intégrabilité
ou sur une égalité fournie de l'extérieur.

La preuve commence par la série positive
\[
b_n=\Lam(n)n^{-2},\qquad
0\le b_n\le 2n^{-3/2}\ (n\ge1),\qquad \sum_n b_n\le6.
\]
Elle construit la convergence uniforme de la série $D$ et sa borne.
Les identités d'inversion Gamma, l'holomorphie sur le demi-plan droit et
les majorants intégrables du noyau proviennent de dépendances indépendantes
déjà compilées. Pour chaque $n>0$, le transport de branche
$(nw)^{-s}=n^{-s}w^{-s}$ est payé avec $n$ réel positif. Enfin la
sommabilité des intégrales des normes autorise l'échange réel entre la
série en $n$ et l'intégrale en $t$. Les termes $\Lam(p^k)$ restent présents.

Sur la droite $\Re s=2$, l'identification classique de cette série avec
$-\zeta'(s)/\zeta(s)$ appartient au raccord zêta distinct. Le théorème
\ref{thm:mellin} utilise directement la vraie série $D$ ; son crédit ne
vaut pas validation de toutes les formules explicites en zéros, de tous
les déplacements de contour ni d'un évaluateur uniforme de $\zeta$.

L'identité \eqref{eq:mellin} transpose exactement le couplage en une
représentation continue. Elle ne fournit encore aucun signe favorable du
coefficient \eqref{eq:circle}. La domination absolue assure l'existence
des intégrales et les échanges ; une annulation signée nouvelle demeure
une obligation séparée.

\clearpage
\section{Troncatures signées et géométrie du demi-plan droit}\label{sec:tails}
On conserve les définitions concrètes des paquets de troncature. Pour
$\Re w>0$, soit
\begin{equation}\label{eq:Cdelta}
\eta(w)=\frac{\pi/2+|\Arg w|}{2},\quad
\delta(w)=\frac{\pi/2-|\Arg w|}{2}>0,\quad
C(w)=|w|^{-2}\sec^2\eta(w).
\end{equation}
Le noyau Gamma dispose de majorants exponentiels construits sur ce domaine.
Le lot 30, indépendant et compilé, paie les queues non pondérées et leurs
deux orientations. Le bridge du lot 32 relie l'inversion Gamma complète
à la troncature correspondante. Ces crédits n'incluent pas automatiquement
la nouvelle queue pondérée par $D$.

\source Le paquet $E_\Lam$ définit, pour $H\ge0$,
\begin{align}
P_H(w)&=\frac1{2\pi}\int_{-H}^{H}D(t)K(w,t)\,dt,\\
T_+(w,H)&=\int_H^\infty D(t)K(w,t)\,dt,\qquad
T_-(w,H)=\int_H^\infty D(-t)K(w,-t)\,dt.
\end{align}
Il vise les énoncés exacts
\begin{equation}\label{eq:lambda-tail}
P(w)-P_H(w)=\frac{T_-(w,H)+T_+(w,H)}{2\pi},\qquad
|P(w)-P_H(w)|\le
\frac{6C(w)e^{-\delta(w)H}}{\pi\delta(w)}.
\end{equation}
Le facteur 6 vient du théorème \ref{thm:mellin}. La réflexion concerne les
deux facteurs $D$ et $K$. La source comporte des preuves concrètes
d'intégrabilité et de queue ; aucune prémisse égale à
\eqref{eq:lambda-tail} n'est ajoutée.

Le lot 36 échoue sur une réécriture de la réflexion à la ligne 118:6.
Il n'a produit aucun objet de preuve compilé pour ce module ; la géométrie
suivante n'a pas été invoquée. La réparation est conservée dans un nouveau
paquet SOURCE02, sans nouvelle compilation à la clôture de cette étape.
Le diagnostic est un raccord d'élaboration, sans réfutation analytique
ni preuve d'obstruction de parité.

\source Le paquet Geometry25 construit, pour $a>0$,
\begin{equation}\label{eq:geometry}
C(a-i\theta)=\frac2{a^2}
\left(1+\frac{|\theta|}{\sqrt{a^2+\theta^2}}\right)\le\frac4{a^2}.
\end{equation}
Pour $|\theta|\le\pi$, il vise également
\begin{equation}\label{eq:d}
\delta(a-i\theta)\ge d(a):=\tfrac12\arctan(a/\pi)>0.
\end{equation}
Les branches, quadrants, signes, racines carrées et continuités sont
explicitement rédigés. Les rayons sont fermés et continus sur leurs
domaines : $\Re w>0$ et $H\in\R$ pour leur expression, $H\ge0$ pour la
borne de queue. Cette continuité d'une expression ne vaut pas une preuve
supplémentaire de continuité d'une intégrale à frontière mobile.

\clearpage
\section{Enveloppe de coefficient et raccord scalaire restant}\label{sec:radius}
À partir de \eqref{eq:lambda-tail}--\eqref{eq:d}, le rayon uniforme proposé est
\begin{equation}\label{eq:eps}
\varepsilon(a,H)=\frac{24e^{-d(a)H}}{\pi a^2d(a)},\qquad a>0,\ H\ge0.
\end{equation}
La source du coefficient définit la projection centrée de la troncature
\[
C_{N,H}=\frac{e^{aN}}{2\pi}\int_{-\pi}^{\pi}
P_H(a-i\theta)^2e^{-iN\theta}\,d\theta
\]
et le rayon fermé
\begin{equation}\label{eq:E-mellin}
\mathcal E_N(a,H)=e^{aN}\bigl(2U(a)\varepsilon(a,H)
+\varepsilon(a,H)^2\bigr).
\end{equation}
\source Les 34 déclarations de ce paquet construisent le transport du
cercle centré, la continuité et l'intégrabilité angulaires à $H$ fixé,
puis visent
\begin{equation}\label{eq:coef-error}
|C_N-C_{N,H}|\le\mathcal E_N(a,H).
\end{equation}
L'inégalité élémentaire employée est
$|x^2-y^2|\le2|x||x-y|+|x-y|^2$. Avec $|T_a|\le U(a)$, elle explique
les deux termes du budget. La continuité jointe du rayon est rédigée pour
$a>0$ et $H\in\R$ ; son interprétation comme erreur de troncature exige
$H\ge0$. Le théorème n'est pas compilé à la clôture.

La vraie série $T_a$ est périodique. La puissance complexe
$(a-i\theta)^{-2-it}$, prise isolément sur sa branche principale, ne l'est
pas. La source ne postule pas la périodicité de $P_H$. Le déplacement de
l'intervalle de projection concerne la trace complète, dont la périodicité
provient de la série ; la troncature reçoit ensuite son propre budget.

\source Le paquet scalaire de 26 déclarations définit les mêmes fonctions
$U$, $d$, $\varepsilon$ et $\mathcal E_N$, et rédige une chaîne réelle au
point
\begin{equation}\label{eq:point}
N=10^8,\qquad a=10^{-8}=1/N,\qquad H=10^{11},\qquad\tau=10^{-6}.
\end{equation}
Il vise notamment $d(a)>1/(8N)$, une borne de $U$ par l'identité en
$\sinh$, $e^{125}>10^{50}$, puis
$\mathcal E_N(a,H)\le576/10^{10}<\tau$.
Il n'est pas compilé. Un bridge distinct de huit théorèmes rédige les
égalités des deux définitions de rayon et l'instanciation de
\eqref{eq:coef-error} par ce résultat scalaire. Ses deux imports sont
prospectifs ; le bridge n'est pas compilé non plus.

Ces sources traitent des obligations théoriques précises sans établir un
coefficient numérique. Avant leur compilation indépendante, aucune chaîne
Lean complète n'autorise à promouvoir \eqref{eq:coef-error} et son
instanciation au point \eqref{eq:point} en résultat formel acquis.

\clearpage
\section{Le test rationnel effectivement terminé et sa portée}\label{sec:scalar-test}
Le test scalaire observé emploie uniquement des fractions exactes. Au point
\eqref{eq:point}, il définit
\[
S_{150}(125)=\sum_{k=0}^{150}\frac{125^k}{k!},\qquad
q=\frac1{S_{150}(125)},\qquad
\mathcal B=3\bigl(128N^5q+4096N^6q^2\bigr).
\]
\textbf{Calcul observé : PASS rationnel.} Une invocation réelle, sans
réessai, vérifie exactement
\begin{equation}\label{eq:B-test}
\mathcal B<10^{-6}.
\end{equation}
Les contrôles conservés incluent
$S_6(5)>100$ et
$S_{150}(125)\ge S_6(5)^{25}>10^{50}$.
Le résultat est lié au reçu de l'acteur et à l'observation physique ROOT
répertoriés en annexe. Il n'emploie aucune approximation flottante.

La comparaison \eqref{eq:B-test} est un résultat scalaire. Le programme
n'intègre pas $D$, ne calcule ni $P_H$ ni $C_{N,H}$, et n'évalue pas
$C_N$. Il ne vérifie pas la conformité d'un catalogue de nombres premiers,
la transformée native ou une identité de Goldbach. Le raccord entre
$\mathcal B$ et le rayon réel \eqref{eq:E-mellin} utilise les inégalités
analytiques rédigées dans le paquet scalaire et le bridge de la section
\ref{sec:radius}. Leur statut reste SOURCE.

\paper Le raccord analytique au même rayon s'écrit explicitement comme
suit. Les inégalités réelles des connecteurs donnent $3<\pi<7/2$,
$d(1/N)>1/(8N)$ et
$U(1/N)=1/(4\sinh^2(1/(2N)))\le N^2$.
Ainsi $dH>125$ au point fixé. La série de Taylor positive assure
$S_{150}(125)\le e^{125}$, donc
$e^{-dH}\le e^{-125}\le q$. L'expression \eqref{eq:eps} vérifie alors
\[
\varepsilon(1/N,H)\le
\frac{24}{3}N^2(8N)q=64N^3q.
\]
Comme $aN=1$ et $e<3$, \eqref{eq:E-mellin} implique
\[
\mathcal E_N(1/N,H)\le
3\bigl(2N^2\cdot64N^3q+(64N^3q)^2\bigr)=\mathcal B<\tau.
\]
Cette chaîne explique exactement la portée analytique proposée du test.
Elle reste PAPER/SOURCE jusqu'à compilation des connecteurs et du bridge
au coefficient ; le calcul rationnel seul ne vérifie pas ces dépendances.

\subsection*{Faisabilité d'une évaluation Mellin littérale}
\paper La géométrie serrée donne le rayon uniforme
\[
\overline C(a)=\frac2{a^2}\left(1+
\frac{\pi}{\sqrt{a^2+\pi^2}}\right)<\frac4{a^2},\qquad
\overline\varepsilon=
\frac{6\overline C(a)e^{-d(a)H}}{\pi d(a)}.
\]
La note de coût étudie le seuil du majorant
$e^{aN}(2U\overline\varepsilon+\overline\varepsilon^2)$.
Pour les paramètres \eqref{eq:point}, elle place ce seuil entre
$5\cdot10^{10}$ et $10^{11}$ par des inégalités fermées.
Cette estimation porte sur une condition suffisante issue d'un majorant ;
elle n'est pas une minoration de l'erreur réelle et n'impose pas cette
hauteur à toutes les méthodes possibles.

La même note paie sur papier une queue littérale de $D$ par
$(\log M+1)/M$ et une quadrature angulaire par un majorant de dérivée.
Des choix extrêmement grossiers tels que $M=10^{52}$ et $J=10^{67}$
illustrent la fermeture symbolique de ces budgets. Ils ne décrivent ni
un programme disponible ni une exécution achevable dans la fenêtre
de ressources de cette étape. La certification dirigée des primitives
sur tous les n\oe{}uds et la quadrature verticale demeurent ouvertes.
Le succès scalaire \eqref{eq:B-test} n'est donc pas une mesure de
performance de l'évaluation globale de Mellin.

\clearpage
\section{Contrat numérique falsifiable et arrêt natif sans verdict}\label{sec:numeric}
Le contrat fini de la section \ref{sec:crt} donne une cible falsifiable :
produire $I_N$ depuis les coefficients canoniques, reconstruire le même
entier par les cinq résidus, et vérifier toutes les gardes. Alors
\eqref{eq:E-log} fournit l'intervalle
\[
C_N\in\left[I_N/S^2-E_{\log},\ I_N/S^2+E_{\log}\right].
\]
Le rayon n'est pas choisi après observation ; sa construction réelle et
sa garde sont déjà compilées. La validité d'un résultat natif exige encore
la conformité des enregistrements à $A32$, la couverture complète du
catalogue et les vérifications du calcul entier.

\paper Pour $K=2^{27}$, une transformée radix 2 possède
$K\cdot27/2=1\,811\,939\,328$ papillons par transformée. Les comptes du
contrat donnent des dizaines de milliards de produits modulaires pour
les cinq modules et leur contrôle indépendant. Un tableau de $K$ mots
de 32 bits occupe 536\,870\,912 octets ; les points en 64 bits sur
$0,\ldots,N$ occupent 800\,000\,008 octets, et les facteurs en 32 bits
400\,000\,004 octets. Ces comptes de charges utiles ne sont pas une
mesure de mémoire totale ni une promesse de durée.

\textbf{Calcul observé : STOP sans verdict.} Le consumer05 réutilise
trois payloads gelés du producteur04 dans des copies identiques, et une
image checker issue du BUILD04. Les payloads demeurent non validés tant
que le checker n'a pas terminé son contrat. La tentative crée un seul
checker, aucun producteur et aucun compilateur, sans réessai.
Le budget parent est de 3\,600 secondes ; le checker dispose d'une
échéance à 3\,300 secondes, avec une réserve de 300 secondes pour les
journaux et les contrôles de conservation.

Le checker s'arrête avec le code 238 et le statut
\texttt{HARD\_WALL\_NO\_VERDICT}. Le parent termine avec code 1 après
environ 3\,312 secondes, soit 55 minutes. Aucun coefficient n'est présent.
La fermeture documente la quiescence et le POST réellement produit ;
l'observation ROOT vérifie la conservation de 6\,576 entrées, 109 paires
de captures, trois paires de payloads et 3\,089 fichiers d'archives.
L'arrêt constitue un échec de ressources sans verdict mathématique,
et non un contre-exemple à l'identité ou à la précision construite.

Les ressources OS distinguent une limite de mémoire engagée du Job de
2 Gio et un suivi RSS qui n'est pas une limite RSS imposée par l'OS.
La confiance opérationnelle Windows/GCC/natif reste déclarée, à portée
bornée ; une fermeture universelle du chargeur n'est pas affirmée.
Aucune nouvelle tentative ni extrapolation de débit n'est intégrée
dans cette clôture.

\clearpage
\section{Raccord au bilan : la dette signée demeure}\label{sec:dn}
Le coefficient $C_N$ contient des puissances premières propres.
Définissons directement
\[
P_{\mathrm{pr}}(n)=\begin{cases}\log n,&n\text{ premier},\\0,&\text{sinon},\end{cases}
\qquad V_{\mathrm{pp}}(n)=\Lam(n)-P_{\mathrm{pr}}(n).
\]
La notation $V_{\mathrm{pp}}$ évite de confondre ce poids avec l'amplitude
$U(a)$. Dans l'audit du ledger, le coefficient de paires premières et la
correction sont
\begin{align}
R_{\mathrm{sf}}&=\sum_{n=0}^N P_{\mathrm{pr}}(n)P_{\mathrm{pr}}(N-n),\\
Q_N&=C_N-R_{\mathrm{sf}}\nonumber\\
&=\sum_{n=0}^N\bigl(2P_{\mathrm{pr}}(n)V_{\mathrm{pp}}(N-n)
+V_{\mathrm{pp}}(n)V_{\mathrm{pp}}(N-n)\bigr)\ge0.
\label{eq:PP}
\end{align}
\paper Le majorant direct $Q_N\le2\sqrt N\,u^2$ pour $N\ge2$
est conservé dans cette note. Il n'est pas crédité ici comme un nouveau
théorème Lean. Ce terme global n'est pas interchangeable avec les termes
de puissances premières dont le bilan garde ses propres supports.

Le pont couvert de la monographie, avec ses conditions et exactement
le même support, conserve
\begin{equation}\label{eq:ledger}
D_N+\varepsilon_{\mathrm{tr}}
=R_{\mathrm{ref}}-R_{\mathrm{sf}}+2\max(e,0).
\end{equation}
Son raccord au coefficient continu est donc
\begin{equation}\label{eq:gap}
D_N=R_{\mathrm{ref}}-C_N+Q_N+2\max(e,0)-\varepsilon_{\mathrm{tr}}.
\end{equation}
Le ledger conserve une fois chacun
$B_{\mathrm{prime}}^{(a_{\mathrm{cut}})}$,
$B_{\mathrm{pp}}^{(a_{\mathrm{cut}})}$,
$P_{\mathrm{band},\ge2}$, $Z_{\mathrm{face},\ge2}$, $I_\alpha$ et
$2\max(e,0)$. La présente transposition ne redéveloppe pas ces objets
par les méthodes arithmétiques abandonnées dans la boucle 22.

La connaissance d'un intervalle étroit pour $C_N$ ne donne pas le signe
de $R_{\mathrm{ref}}-C_N$, ni le contrôle de $e$, des fronts et des coins.
Pour établir \eqref{eq:target}, il reste à démontrer une estimation
signée compatible avec \eqref{eq:gap}, en conservant ces charges et les
conditions du pont. Poser cette estimation comme une prémisse égale à
la cible ne constituerait pas un contournement du mur de la parité.

Les identités continues établissent des représentations exactes et les
budgets décrivent une erreur de troncature ou de quantification. Aucune
annulation globale de la vraie contribution spectrale transportée au
ledger n'a été prouvée. Cette dette reste ouverte même si un futur
calcul fini produit $C_N$ au point $10^8$. Elle doit être traitée au
seuil \eqref{eq:onset} et avec la couverture finie requise par le cadre.

\clearpage
\section{Bords, zéros et invariance : pistes à portée précise}\label{sec:strategy}
Le pivot continu conserve la phase avant de prendre une norme. Pour
$a>0$, $N\ge1$ et $q\in\C$, on considère
\[
J_{a,N}(q)=\frac1{2\pi}\int_{-\pi}^{\pi}
(a-i\theta)^{-q}e^{-iN\theta}\,d\theta.
\]
\paid Le lot 25 formalise une intégration par parties et le vrai bord
sur la branche principale :
\begin{equation}\label{eq:bord}
J_{a,N}(q)=\frac{i(-1)^N}{2\pi N}
\bigl((a-i\pi)^{-q}-(a+i\pi)^{-q}\bigr)
+\frac qN J_{a,N}(q+1).
\end{equation}
Le bord ne s'annule pas point par point. Ce module de 30 déclarations
ne donne aucun signe favorable à la projection spectrale complète.
La périodicité de la trace après inversion ne rend pas périodique
chaque noyau tronqué.

\paper La stratégie de formule explicite signée localise une fonction
test lisse près de $x+y=N$, puis vise la décomposition
\[
C_N=\mathcal B_N-2\mathcal L_N+\mathcal Z_N.
\]
Les trois termes désignent respectivement fond, couplage au canal des
zéros, et contribution à deux zéros dans la formule explicite envisagée.
L'identité complète, les corrections archimédiennes, les multiplicités,
les échanges et les queues de zéros demeurent à payer dans ce canal.
Les revues conservées constatent la cohérence sous ces dettes ; elles
ne créditent ni un signe de $\mathcal Z_N-2\mathcal L_N$, ni une preuve
Lean de la formule globale.

L'organisation en orbites conjugaison--réflexion conserve les zéros hors
de la droite critique. Une orbite de quatre points engendre une densité
de type
\[
4m x^{-1/2}\cosh\bigl((\beta-1/2)\log x\bigr)
\cos(\gamma\log x),
\]
et une orbite critique de deux points la densité
$2m x^{-1/2}\cos(\gamma\log x)$, avec multiplicité $m$.
Les stabilisateurs empêchent de compter deux fois une orbite dégénérée.
La phase cosinus peut changer de signe : une organisation correcte des
orbites n'est pas une positivité automatique du coefficient additif.

La revue du noyau canonique montre une direction négative dans une
restriction de dimension deux ; elle exclut l'argument de positivité
naïve examiné. Sa portée est ce noyau précis. Elle n'est ni une
impossibilité universelle de toute approche opératorielle, ni une
réfutation de Goldbach. De même, une modularité ou une invariance
exacte proposée devra être démontrée pour les vrais objets et devra
isoler le terme signé requis dans \eqref{eq:gap}. Aucun tel raccord
suffisant n'est annoncé comme imminent dans cette clôture.

\clearpage
\section{Fejér continu : reconstruction et signe restant}\label{sec:fejer}
La nouvelle note spectrale part de la vraie série $D$ du lot 35, avec
$c_n=\Lam(n)/n^2$ pour $n\ge2$ et $c_0=c_1=0$.
La masse $\sum c_n\le6$ rend $D$ de type positif : pour des familles
finies $(z_j,t_j)$,
\[
\sum_{j,k}\overline z_jz_kD(t_k-t_j)
=\sum_{n\ge2}c_n\left|\sum_k z_ke^{-it_k\log n}\right|^2\ge0.
\]
\paper Pour $T>0$ et $m\ge1$ entier, la moyenne triangulaire
\begin{equation}\label{eq:fejer-cont}
\begin{split}
A_T(m)&=\frac1T\int_{-T}^{T}\left(1-\frac{|t|}{T}\right)
D(t)e^{it\log m}\,dt\\
&=\sum_{n\ge2}c_n\operatorname{sinc}^2\!\left(\frac T2\log\frac mn\right)
\end{split}
\end{equation}
est réelle et positive. Ici $\operatorname{sinc}(x)=\sin(x)/x$ si
$x\ne0$ et $\operatorname{sinc}(0)=1$. L'échange sur le compact est
payé par la sommabilité ; le triangle est l'autocorrélation d'une
indicatrice d'intervalle. Les fréquences distinctes vérifient
$|\log(m/n)|\ge\log(1+1/m)$. La note dérive donc le rayon fermé
\begin{equation}\label{eq:fejer-radius}
0\le A_T(m)-c_m\le
\frac{24}{T^2\log^2(1+1/m)},\qquad T>0.
\end{equation}
Le cas $m=1$ a $c_1=0$ et conserve la même borne. Le rayon est continu
en $T>0$ et tend vers zéro à $m$ fixé. Ni cette identité ni ses
connecteurs ne sont compilés dans cette étape.

Ce mécanisme reconstruit des amplitudes depuis la donnée verticale vraie,
puis peut les insérer dans les blocs réfléchis $n,N-n$. Il ne crée pas
de masse commune à une paire. La note conserve un contre-modèle borné :
le seul secteur des puissances de 2 est de type positif et satisfait la
même reconstruction, alors que son coefficient additif à $N=14$ est nul.
Ce constat exclut la déduction d'un minorant strict uniforme depuis cette
seule positivité. Il ne remplace pas la vraie zêta et n'exclut pas une
identité supplémentaire spécifique à sa donnée complète.

Un évaluateur futur devrait produire les intervalles des vraies
intégrales $A_T(m)$, avec les primitives, positions et quadratures
contrôlées. \eqref{eq:fejer-radius} ne paie aucune de ces erreurs
d'évaluation. Le coût augmente avec $m$, et aucune faisabilité à $10^8$
n'est annoncée. Les identités anciennes de grilles divisibles et leurs
covariances demeurent dans V2, avec leurs conventions de supports ; la
présente piste concerne la variable spectrale continue $t$.

La recherche est gelée. Une reprise doit établir un raccord signé
supplémentaire au ledger \eqref{eq:gap}, plutôt que rebaptiser cette
reconstruction comme annulation de parité. Aucun mécanisme suffisant
pour \eqref{eq:target} ni découverte imminente n'est revendiqué.

\clearpage
\section{Inventaire formel arrêté et journal des limites}\label{sec:inventory}
L'observation indépendante du lot 35 fixe l'inventaire officiel à
\textbf{88 modules / 1\,488 déclarations auxiliaires, définitions incluses}.
Le lot 36 n'ajoute aucun crédit. Le tableau suivant isole les éléments
utilisés dans cette synthèse ; il n'est pas une somme exhaustive de
l'inventaire historique.

\begingroup\small\setlength{\tabcolsep}{3pt}
\noindent
\begin{tabularx}{\textwidth}{@{}p{12mm}Yp{22mm}p{25mm}@{}}
\toprule
Lot & Objet & Déclarations & Statut\\
\midrule
11 & Enveloppe géométrique de la vraie trace thermique & 42 & Row PASS ; lot FAILED\\
13 & Identité continue d'extraction & 20 & Lean PASS\\
15 & Orthogonalité et projecteur discret complexe & 29 & Lean PASS\\
17 & Logarithmes canoniques et enveloppe & 52 & Lean PASS\\
19 & Projecteur générique de corps fini & 24 & Lean PASS\\
22 & Cinq racines concrètes et projection $A32$ & 28 & Lean PASS\\
25 & Intégration par parties et terme de bord & 30 & Lean PASS\\
26/27 & Inversion Gamma locale et holomorphe & 22/15 & Row PASS26 ; PASS27\\
30/32 & Queues Gamma et bridge non pondéré & 22/2 & Lean PASS\\
35 & Mellin de la vraie série $\Lam$ & 30 & Lean PASS\\
36 & Queue pondérée $E_\Lam$ ; géométrie & 16 ; 25 & FAIL ; non invoqué\\
\bottomrule
\end{tabularx}
\endgroup

Les lots 11 et 26 ont un statut global FAILED : chacun conserve une
ligne de module entier PASS (respectivement Envelope42 et Local22),
tandis qu'un autre module du même lot échoue. Leurs crédits partiels
sont attestés par les adjudications et observations ; les modules
défaillants ne sont pas comptés comme acquis. Le lot 13 et le lot 27
valident ultérieurement des sources corrigées distinctes.

Les sources du coefficient (34), des connecteurs scalaires (26) et du
bridge coefficient--scalaire (8) n'ont pas été compilées. Leurs nombres
de déclarations sont des catalogues de sources ; ils ne s'ajoutent pas
à 1\,488. Les anciens commentaires \texttt{SOURCE ONLY} de certains
fichiers devenus PASS restent dans les captures historiques ; le statut
actuel vient de l'adjudication et du reçu, pas d'une réécriture rétroactive.

Les lots de correction ont conservé leurs échecs : le lot 16 présente
des raccords de série et de conversions rationnelles, réparés et compilés
dans le lot 17 ; les lots 31 et 33 présentent des raccords d'instances
ou de continuité, réparés dans la source effectivement compilée au lot
35. Le lot 36 échoue sur la réflexion du produit de deux fonctions.
Ces diagnostics ne sont pas des preuves du mur de la parité. À
l'inverse, un PASS auxiliaire ne vaut pas victoire sur cette obstruction.

La vérification est une propriété du fichier exact, du catalogue exact
et de ses dépendances. Les contrôles conservent les anciens journaux,
captures, sources et archives ; ils distinguent lecture intégrale d'un
texte, lecture ciblée, inventaire métadonnées et hachage des octets.
Le fichier compagnon \path{evidence_v3.json} précise cette frontière
pour les principales affirmations de ce manuscrit.

\clearpage
\appendix
\section{Provenance compacte et reprise}\label{app:provenance}
Tous les chemins ci-dessous sont relatifs à la racine du corpus
\path{Goldbach_Parity_20261002}. Les préfixes SHA-256 servent à la lecture ;
les empreintes complètes, tailles, statuts et scopes de lecture sont
conservés dans \path{synthesis_v3/evidence_v3.json}. Les chemins de sources
ne remplacent pas les reçus d'exécution.
Le corpus publié est accessible dans le dépôt
\href{https://github.com/seisner1967-hash/goldbach-horizon}{\texttt{goldbach-horizon}},
avec le registre compagnon de la V3.

\begingroup\small\setlength{\tabcolsep}{3pt}
\noindent
\begin{tabularx}{\textwidth}{@{}p{20mm}Yp{23mm}@{}}
\toprule
Pièce & Chemin relatif & Préfixe SHA\\
\midrule
Mellin PASS & \path{round22/judge5/batch35/sources/ComplexGammaMellinLambda22.lean} & \texttt{4a94415678e4}\\
Reçu 35 & \path{round22/judge5/batch35/batch35_attempt01/receipt.json} & \texttt{3cffa331b392}\\
Observation 35 & \path{.arbor/sessions/parity/.coordinator/messages/round22_judge_batch35_closed_observation.json} & \texttt{dc5985976051}\\
Queue réparée & \path{round22/role4/complex_gamma_mellin_lambda_tail_source02/ComplexGammaMellinLambdaTail22.lean} & Voir JSON\\
Géométrie SOURCE & \path{round22/role4/complex_gamma_circle_geometry_source01/ComplexGammaCircleGeometry22.lean} & \texttt{2322915060c6}\\
Coefficient SOURCE & \path{round22/role4/lambda_circle_truncation_source01/LambdaCircleTruncationEnvelope22.lean} & \texttt{bf9b8257760a}\\
Scalaire SOURCE & \path{round22/role4/scalar_radius_real_connectors_source01/ScalarRadiusRealConnectors22.lean} & \texttt{6ba4b8d4cca0}\\
Bridge SOURCE & \path{round22/role3/lambda_circle_scalar_bridge_source01/LambdaCircleScalarRadiusBridge22.lean} & \texttt{cfb0d2cfd238}\\
Point rationnel & \path{round22/role6/mellin_circle_radius_point_source01/actual_radius_point01_attempt01/receipt.json} & \texttt{d46e3ffeb6d0}\\
Ledger & \path{round22/role4/continuous_dn_gap_paper01/dn_gap_audit22.md} & Voir JSON\\
\bottomrule
\end{tabularx}
\endgroup

\subsection*{Contrat autonome de remise}
La pièce \path{synthesis_v3/HANDOFF_THEORY_V3.txt} résume les objets, les
paramètres, les acquis et les dettes. Le calcul de $C_N$ à $10^8$ demeure
absent. La reprise native éventuelle exigerait un nouveau contrat de
ressources et un verdict complet ; les sorties
STOP ne peuvent servir d'oracle. Une reprise d'ingénierie analytique
éventuelle devra compiler indépendamment les sources de queue,
géométrie, coefficient et connecteurs, puis payer les primitives et
quadratures réellement choisies. La reprise théorique sur le signe de
$D_N$ peut partir du socle compilé, en conservant ces statuts ouverts.
Une borne signée sur le ledger demeure l'obligation centrale.

\subsection*{Sources documentaires locales}
Cette synthèse utilise la monographie antérieure, les sources Lean gelées,
les adjudications et les revues réellement consultées. Elle ne prétend pas
avoir vérifié de nouvelles références externes ni les citations secondaires
des notes. La bibliographie ci-dessous identifie les documents locaux qui
portent les arguments écrits et leurs limites.
\begin{thebibliography}{9}\small
\bibitem{v2} Durand Serge, \emph{Finite Identities and Explicit Error Budgets
in a Framework for Binary Goldbach}, synthèse V2, 1 octobre 2026.
Titre, auteur et continuité relevés dans la source publique conservée.
\bibitem{mono} Durand Serge, monographie finale et extraction locale
\path{monographie.txt} ; conventions actualisées par le cadre source
\path{Goldbach_Final_20261001/goldbach_synthesis.tex}.
\bibitem{ntt} ROLE4, \emph{Contrat de projection fini},
\path{round22/role4/circle_ntt_paper01/projection_contract22_final.md}.
\bibitem{cost} ROLE4, \emph{Faisabilité Mellin--Lambda},
\path{round22/role4/mellin_lambda_feasibility_paper01/feasibility22.md}.
\bibitem{gap} ROLE4, \emph{Audit du raccord spectral à $D_N$},
\path{round22/role4/continuous_dn_gap_paper01/dn_gap_audit22.md}.
\bibitem{weil} Juge indépendant, \emph{Revue de formule explicite signée},
\path{round22/judge5/signed_weil_source_review01/review22.md}.
\bibitem{orbit} ROLE4, \emph{Revue des orbites de zéros},
\path{round22/role4/zero_orbit_source_review01/review22.md}.
\bibitem{fejer} ROLE1, \emph{Positivité réelle de $D$ et reconstruction
des blocs conservés},
\path{round22/role1/spectral_next_paper01/spectral_next22.md}.
\end{thebibliography}
\end{document}
